% ================================================================
% cover_letter_v4.0.tex — Cover letter for JMST submission
% ================================================================
\documentclass[12pt,a4paper]{article}
\usepackage[margin=2.5cm]{geometry}
\usepackage{amsmath,graphicx,hyperref}

\begin{document}

\begin{flushright}
Zhilei Chen\\
Guangdong Peizheng College\\
Data and Computer Science\\
53 Peizheng Rd., Chini Town, Huadu\\
Guangzhou, Guangdong 510830, China\\
E-mail: 2604513@peizheng.edu.cn\\[6pt]
\today
\end{flushright}

\vspace{0.5cm}

Editor-in-Chief\\
Journal of Mechanical Science and Technology (JMST)\\
Korean Society of Mechanical Engineers

\vspace{0.8cm}

\textbf{Re: Submission of ``Method-Induced Uncertainty in Gear Bending
Fatigue: Factorial Decomposition of Size-and-Surface Factor Sensitivity
Under a Common ISO Stress Baseline''}

\vspace{0.5cm}

Dear Editor,

I am pleased to submit the enclosed manuscript for consideration in the
Journal of Mechanical Science and Technology.

\vspace{0.3cm}
\noindent\textbf{1. Clear Problem Statement}

When three engineers---using different standards (ISO, AGMA, FKM),
different mean-stress corrections, and different surface-finish
assumptions---compute the bending fatigue life of the \emph{same} spur
gear under the \emph{same} load, their predictions span three orders of
magnitude. Prior gear fatigue studies have compared individual standards
or correction methods in isolation (e.g., Bonaiti 2024 for ISO vs.\
AGMA; Dowling 2009 for mean-stress corrections). The present study
extends this tradition by simultaneously decomposing the contributions
of five methodological choices and ranking them by variance explained.

\vspace{0.3cm}
\noindent\textbf{2. Critical Methodology}

A censored-likelihood AFT survival model is applied to a 144-combination
full-factorial design (5 factors $\times$ 2--4 levels, 92 run-outs).
Unlike ordinary ANOVA/Sobol'---which must discard all 92 run-out entries
and produce systematically distorted variance fractions---the AFT
framework retains right-censored observations through the EM algorithm.
Bootstrap (2000 draws), Sobol' sensitivity analysis, Weibull-AFT model
selection ($\Delta\mathrm{AIC}=+29.7$ in favor of log-normal), and full
residual diagnostics (Shapiro-Wilk $p=0.599$, Levene's test) confirm the
model specification. Sensitivity analyses for the $Y_{Sa}$ stress
concentration assumption ($\leq+2.1$~pp bias) and the application factor
$K_A$ (1.0--1.5, rankings stable) validate the robustness of the
findings.

\vspace{0.3cm}
\noindent\textbf{3. Consequential Findings}

Three factors account for 97.9\% of total log-variance: mean-stress
correction (39.1\%), size-and-surface standard (35.3\%), and surface
roughness (23.5\%). The S--N curve source ($-$0.7\%) and cumulative
damage rule (0.3\%, constant-amplitude only) are negligible. This
ranking is stable under bootstrap resampling, load-factor variation, and
alternative stress-concentration treatments. Engineering design
recommendations---including calibration factors between standards and
risk-classification criteria---are derived directly from the statistical
results.

\vspace{0.3cm}
\noindent\textbf{4. Fit to JMST Scope}

The manuscript aligns with JMST's Solid Mechanics \& Design section.
It addresses a practical mechanical engineering problem (gear fatigue
life prediction), employs rigorous statistical methods (censored AFT),
quantifies engineering uncertainty sources in terms accessible to design
engineers, and provides a reusable computational framework applicable to
other machine elements. No physical experiments are required---the
entire analysis uses publicly available S--N data and freely verifiable
standard formulas---making the work fully reproducible.

\vspace{0.3cm}
\noindent\textbf{5. Manuscript Details}

\begin{itemize}
  \item Manuscript type: Original Article
  \item Figures: 7; Tables: 3
  \item Supplementary material: Full reproducibility package (Python
        code, sweep data, AFT decomposition results, sensitivity
        analysis scripts)
  \item The manuscript has been formatted for double-blind review per
        JMST guidelines
\end{itemize}

\vspace{0.3cm}
I confirm that this manuscript has not been published elsewhere and is
not under consideration by another journal. I have approved the
manuscript and agree with its submission to JMST.

\vspace{0.5cm}
\noindent Thank you for considering this submission. I look forward to
the reviewers' feedback.

\vspace{0.8cm}
\noindent Sincerely,\\[12pt]
Zhilei Chen
\end{document}
